\section{Construcción de la bitácora}\label{sec:construccion-bitacora}

Con el propósito de establecer una base teórica sólida que permitiera sustentar la toma de decisiones respecto a la selección de un servicio de directorio, se llevó a cabo una revisión exhaustiva de la literatura relacionada con el tema. Este proceso se desarrolló en diferentes etapas, las cuales se describen a continuación:


\subsection{Planeación}\label{subsec:planeacion}

El objetivo de esta etapa fue definir el propósito general del \ac{SMS}, adoptando la metodología propuesta por \citet{Sepulveda2021}. A su vez, se definieron metas (ver \autoref{tab:metas}), preguntas de investigación (ver \autoref{tab:preguntas}) y métricas (ver \autoref{tab:metricas}). Para ello, se siguió el modelo Objetivo-Pregunta-Métrica (\acs{GQM}).


\subsubsection{Definición de metas para el \texorpdfstring{\ac{SMS}}{SMS}}
\begin{table}[H]
	\centering
	\caption{Definición de metas del \acs{SMS}}
	\label{tab:metas}

	\renewcommand{\arraystretch}{1.2}
	\fontsize{9pt}{10pt}\selectfont

	\begin{tabularx}{\textwidth}{|c|X|}
		\hline
		\textbf{Meta} &
		\textbf{Descripción}
		\\
		\hline
		G1            &
		Identificar trabajos acerca de herramientas tecnológicas relacionadas con servicios de directorio de código abierto/software libre.
		\\
		\hline
		G2            &
		Analizar y clasificar las herramientas tecnológicas relacionadas con servicios de directorio de código abierto según sus características técnicas, arquitectónicas y casos de uso reportados.
		\\
		\hline
	\end{tabularx}
\end{table}



\subsubsection{Definición de preguntas de investigación}\label{subsubsec:rq-def}
\begin{table}[H]
	\centering
	\caption{Definición de preguntas de investigación del \acs{SMS}}
	\label{tab:preguntas}

	\renewcommand{\arraystretch}{1.2}
	\fontsize{9pt}{10pt}\selectfont

	\begin{tabularx}{\textwidth}{|c|c|X|X|}
		\hline
		\textbf{Meta}                                                                                                                                                                                 &
		\textbf{Pregunta}                                                                                                                                                                             &
		\textbf{Descripción}                                                                                                                                                                          &
		\textbf{Motivación}
		\\
		\hline
		G1                                                                                                                                                                                            &
		Q1                                                                                                                                                                                            &
		\textit{¿Qué herramientas tecnológicas de código abierto se relacionan con servicios de directorio?}                                                                                          &
		\textbullet\ Identificar las soluciones de servicios de directorio de código abierto existentes en la literatura

		\textbullet\ Reconocer cuáles tecnologías han sido más estudiadas o utilizadas en distintos contextos.
		\\
		\hline
		G2                                                                                                                                                                                            &
		Q2                                                                                                                                                                                            &
		\textit{¿Cuáles son las características técnicas, arquitectónicas y los contextos de aplicación de las herramientas tecnológicas relacionadas con servicios de directorio de código abierto?} &
		\textbullet\ Reconocer cómo están estructuradas arquitectónica y técnicamente las soluciones de servicios de directorio de código abierto.

		\textbullet\ Identificar las características funcionales y técnicas más relevantes reportadas.

		\textbullet\ Analizar los contextos de aplicación y tendencias de uso de estas soluciones.
		\\
		\hline
	\end{tabularx}
\end{table}



\subsubsection{Definición de métricas}
\begin{table}[H]
	\centering
	\caption{Definición de métricas del \acs{SMS}}
	\label{tab:metricas}

	\renewcommand{\arraystretch}{1.2}
	\fontsize{9pt}{10pt}\selectfont

	\begin{tabularx}{\textwidth}{|c|X|}
		\hline
		\textbf{Métrica} &
		\textbf{Descripción}
		\\
		\hline
		M1               &
		Frecuencia de aparición de cada herramienta tecnológica.
		\\
		\hline
		M2               &
		Frecuencia de características técnicas y arquitectónicas reportadas en las herramientas tecnológicas.
		\\
		\hline
	\end{tabularx}
\end{table}

